% =============================================================================
% Draft mục 3.3.3 — Order Book (OrderStore)
% =============================================================================

\subsection{Order Book (OrderStore)}
\label{sec:order-book}

\texttt{OrderStore} là sổ lệnh (orderbook) của hệ thống, nơi các intent đang active được lưu trữ giữa thời điểm người dùng gửi và thời điểm Matching Engine khớp được chúng. Thiết kế của OrderStore phản ánh một quyết định kiến trúc quan trọng: tách \emph{nguồn sự thật lúc chạy} (runtime source of truth) khỏi \emph{lưu trữ bền} (durable storage), dùng cơ chế write-through cache để vừa giữ độ trễ thấp cho matching tick, vừa đảm bảo dữ liệu không mất khi tiến trình khởi động lại.

\subsubsection{Hai tầng dữ liệu}

\subsubsection*{Tầng trong bộ nhớ (in-memory)}

Tầng này là các cấu trúc JavaScript \texttt{Map} sống trong tiến trình Node của backend, gồm ba cấu trúc chính:

\begin{itemize}
    \item \emph{Primary index} dạng \texttt{Map<orderId, IntentOrder>} cho tra cứu một order theo định danh trong $O(1)$.
    \item \emph{Pair index} dạng \texttt{Map<"src-des", Set<orderId>>} cho tra cứu nhanh tập order đang hoạt động theo cặp chuỗi nguồn--đích trong $O(1)$. Pair index chỉ chứa các order \texttt{QUEUED} hoặc \texttt{PARTIAL}; khi một order chuyển sang \texttt{MATCHED} hoặc \texttt{EXPIRED}, pair index được xoá tham chiếu tương ứng để Matching Engine không phải lọc lại.
    \item \emph{Match log} dạng mảng \texttt{MatchResult[]} append-only, lưu lịch sử tất cả các sự kiện match đã phát sinh.
\end{itemize}

\subsubsection*{Tầng bền vững (Postgres)}

Tầng này là cơ sở dữ liệu quan hệ Postgres chạy trong tiến trình riêng biệt, được truy cập qua Drizzle ORM, với ba bảng tương ứng:

\begin{itemize}
    \item \texttt{intent\_orders}: mỗi dòng là một \texttt{IntentOrder}. Số lượng được lưu dạng \texttt{text} để giữ chính xác tuyệt đối (USDC dùng 6 chữ số thập phân, lớp matching xử lý chuỗi số nguyên bất kỳ).
    \item \texttt{match\_results}: bảng nhật ký append-only của các sự kiện match. Các trường \texttt{cycles} và \texttt{raw\_cycles} được lưu dạng JSONB do hệ thống không cần truy vấn vào nội dung.
    \item \texttt{executions}: lưu trạng thái thực thi mỗi \texttt{matchId}, được cập nhật theo thời gian thực khi Orchestrator chuyển trạng thái \texttt{pending} $\rightarrow$ \texttt{done|error}.
\end{itemize}

\subsubsection{Mô hình write-through}

Mọi thao tác mutating trên OrderStore đều theo trình tự hai bước. Trước tiên, hệ thống cập nhật cấu trúc trong bộ nhớ \emph{đồng bộ} (synchronous) để các thành phần đọc tiếp theo (Matching Engine, các route đọc) nhìn thấy giá trị mới ngay lập tức. Sau đó, hệ thống phát một thao tác ghi tương ứng vào Postgres theo kiểu \emph{fire-and-forget}: tạo Promise, không await, gắn handler \texttt{logDbError} cho trường hợp lỗi.

Mục tiêu thiết kế là không để vòng lặp matching (chu kỳ 1\,s) bị chặn bởi I/O cơ sở dữ liệu. Đánh đổi đi kèm là nếu tiến trình bị crash giữa bước 1 và bước 2, trạng thái trong bộ nhớ và Postgres có thể lệch nhau trong một khoảng ngắn — chấp nhận được ở giai đoạn MVP. Phương án nâng cấp được liệt kê ở mục Future Work là chuyển sang \emph{outbox pattern}, tức append event vào một bảng riêng trong cùng một transaction với việc thay đổi trạng thái, rồi để một consumer riêng đọc bảng đó.

\subsubsection{Hydration on boot}

Vì in-memory không phải durable, nó phải được \emph{rehydrate} (tái xây dựng) từ Postgres mỗi khi tiến trình khởi động. Phương thức \texttt{OrderStore.hydrate()} chạy trước khi \texttt{MatchScheduler} bắt đầu tick, gồm hai bước. Đầu tiên là \texttt{SELECT * FROM intent\_orders} để đọc toàn bộ order, dựng lại primary index và pair index; chỉ các order \texttt{QUEUED} hoặc \texttt{PARTIAL} được đưa vào pair index, vì pair index chỉ đại diện các order còn ứng cử để match. Tiếp theo là \texttt{SELECT * FROM match\_results ORDER BY matched\_at} để đọc nhật ký match và append theo thứ tự thời gian.

Hydration được \emph{await} trong \texttt{start()} của entrypoint backend; nếu hydration thất bại, tiến trình crash thay vì bắt đầu chạy với trạng thái rỗng. Đây là lựa chọn an toàn vì rỗng có thể vô tình bỏ mất các order còn pending hoặc các execution chưa hoàn tất.

\subsubsection{Vai trò trong sơ đồ phụ thuộc}

OrderStore là singleton của module \texttt{matching/store/orderStore} và được chia sẻ bởi hai nhóm consumer chính. Phía API Gateway gọi \texttt{add()} khi nhận \texttt{POST /api/intent} và gọi \texttt{get()} khi nhận các route truy vấn. Phía MatchingService gọi \texttt{getActiveOrders()}, \texttt{getByPair()}, \texttt{update()}, \texttt{removeFromPairIndex()}, \texttt{addMatchResult()}, \texttt{expireStale()}.

Nói cách khác, OrderStore đóng vai trò bus dữ liệu nội bộ: \emph{mọi thay đổi trạng thái orderbook đi qua nó}, và mọi việc ghi xuống Postgres do \emph{một mình OrderStore} chịu trách nhiệm. Nhờ đó các thành phần khác (MatchingService) không cần biết về sự tồn tại của Postgres — một dạng dependency inversion đơn giản.
